\documentclass[conference]{IEEEtran}

\usepackage[utf8]{inputenc}
\usepackage[T1]{fontenc}
\usepackage[brazil]{babel}
\usepackage{amsmath}
\usepackage{booktabs}
\usepackage{graphicx}
\usepackage{microtype}
\usepackage{xcolor}
\usepackage{url}

\graphicspath{{../results/main/figures/}{../results/figures/}}

\title{O Fine-Tuning Altera Apenas a Acurácia ou Também as Regiões Utilizadas por uma CNN?}

\author{\IEEEauthorblockN{Eduardo Jesus Dal Pizzol -- 251008114\\
Matheus Lemes Amaral -- 241031852}
\IEEEauthorblockA{Universidade de Brasília -- UnB\\
FGA0261 -- Inteligência Artificial II}}

\begin{document}
\maketitle

\begin{abstract}
O aprendizado por transferência é amplamente empregado em classificação de imagens, mas sua avaliação costuma se limitar ao desempenho preditivo. Este trabalho investiga se o ajuste completo de uma rede convolucional também modifica as regiões da imagem usadas em suas decisões. Uma ResNet-18 pré-treinada no ImageNet é avaliada no Oxford-IIIT Pet sob duas estratégias: extração de características com o backbone congelado e ajuste fino de todos os parâmetros. Além de acurácia, macro-F1 e calibração, mapas Grad-CAM são comparados com as máscaras de segmentação dos animais por energia no primeiro plano, pointing game e interseção sobre união. O protocolo permite distinguir ganho de desempenho de mudança efetiva na localização da evidência visual e identificar possível dependência do fundo. Os resultados são agregados sobre três sementes e analisados de forma pareada.
\end{abstract}

\begin{IEEEkeywords}
redes neurais convolucionais, transferência de aprendizado, fine-tuning, Grad-CAM, explicabilidade
\end{IEEEkeywords}

\section{Introdução}

Redes neurais convolucionais (CNNs) transformaram a classificação visual ao aprender representações diretamente dos pixels \cite{krizhevsky2012imagenet}. Arquiteturas residuais permitiram treinar modelos mais profundos e tornaram-se backbones comuns em aprendizado por transferência \cite{he2016deep}. Nesse cenário, um modelo pré-treinado em grande escala é adaptado para uma tarefa com menos dados, congelando o extrator de características ou ajustando seus parâmetros.

Embora modelos com melhor desempenho no ImageNet geralmente apresentem melhor transferência, essa relação depende da tarefa e do procedimento de adaptação \cite{kornblith2019better}. Comparações baseadas somente em acurácia não mostram se o modelo aprendeu evidências semanticamente adequadas. Uma CNN pode melhorar a classificação explorando correlações do fundo em vez das características do objeto.

O Grad-CAM produz mapas de importância espacial a partir dos gradientes da classe em relação às ativações convolucionais \cite{selvaraju2017gradcam}. Entretanto, mapas visualmente plausíveis não são evidência suficiente de explicabilidade, pois alguns métodos de saliência podem ser pouco sensíveis aos parâmetros aprendidos \cite{adebayo2018sanity}. Por isso, este trabalho combina inspeção qualitativa e métricas calculadas contra máscaras de segmentação.

São investigadas duas perguntas: \textbf{RQ1}: o fine-tuning completo melhora a classificação em relação ao backbone congelado? \textbf{RQ2}: ele altera a localização da evidência visual e seu alinhamento com o animal? A contribuição é um protocolo reproduzível que relaciona desempenho, calibração e atenção espacial em pares de modelos com a mesma arquitetura e inicialização pré-treinada.

\section{Materiais e Métodos}

\subsection{Dados e particionamento}

É utilizado o Oxford-IIIT Pet, composto por imagens de 37 raças de cães e gatos, acompanhado de rótulos de classe e máscaras de segmentação \cite{parkhi2012cats}. O conjunto oficial de treino e validação é particionado de forma estratificada em 80\% para treinamento e 20\% para validação; o teste oficial permanece isolado. As imagens são redimensionadas e recortadas para $224\times224$ pixels, com normalização do ImageNet. No treinamento são aplicados recorte aleatório, espelhamento horizontal e pequenas variações de cor.

\subsection{Estratégias de transferência}

Os experimentos usam uma ResNet-18 com pesos do ImageNet e uma nova camada linear para 37 classes. Na estratégia \emph{frozen}, o backbone permanece congelado e apenas o classificador é otimizado. Na estratégia \emph{full}, todos os parâmetros são ajustados, usando taxa de aprendizado menor no backbone do que no classificador. AdamW, suavização de rótulos e decaimento cosseno da taxa de aprendizado são usados em ambas as condições.

Cada condição é executada com as sementes 42, 123 e 2026. O melhor checkpoint é escolhido pelo macro-F1 de validação. São reportados acurácia, macro-F1, balanced accuracy e erro esperado de calibração com dez intervalos. Essa combinação evita que classes mais fáceis ou frequentes dominem a análise.

\subsection{Grad-CAM e métricas espaciais}

Para uma classe $c$, o peso do mapa de características $A^k$ é obtido pela média espacial dos gradientes:
\begin{equation}
\alpha_k^c=\frac{1}{Z}\sum_i\sum_j\frac{\partial y^c}{\partial A_{ij}^k}.
\end{equation}
O mapa Grad-CAM é então calculado por
\begin{equation}
L_{\mathrm{GC}}^c=\operatorname{ReLU}\left(\sum_k\alpha_k^c A^k\right),
\end{equation}
e redimensionado para a resolução de entrada. O mapa é gerado para a classe predita por cada modelo na última convolução do bloco residual final.

Três métricas avaliam seu alinhamento com a máscara $M$ do animal. A \emph{energia no primeiro plano} é a fração da intensidade do mapa contida em $M$; o \emph{pointing game} verifica se o máximo do mapa está no animal; e a IoU de saliência compara $M$ com os 20\% pixels de maior ativação. Para cada imagem também são calculadas a similaridade cosseno e a IoU entre as regiões mais salientes dos dois modelos.

As diferenças de métricas espaciais são calculadas por imagem, reduzindo a variabilidade causada pela composição do teste. São apresentadas médias e desvios-padrão entre as três sementes, acompanhados de exemplos de acertos, erros e mudanças de predição.

\subsection{Implementação e reprodutibilidade}
O protocolo foi implementado em um notebook independente para execução no Google Colab, sem dependência do arquivo externo de experimento. A função de carregamento baixa o Oxford-IIIT Pet, constrói as partições estratificadas e aplica as transformações de treinamento e validação descritas anteriormente. Para cada semente, a função de execução instancia duas cópias da ResNet-18, treina primeiro a condição \emph{frozen} e depois a condição \emph{full}, selecionando o checkpoint de maior macro-F1 na validação. Em seguida, o mesmo conjunto de teste é usado para calcular as métricas preditivas e gerar os mapas Grad-CAM.

O código mantém a comparação pareada: os modelos de cada semente recebem a mesma divisão de dados e a mesma inicialização pseudoaleatória, enquanto diferem apenas na política de atualização dos parâmetros. A rotina de análise percorre as imagens de teste, registra a classe predita, a confiança e as três métricas espaciais, e salva os resultados por imagem em arquivos CSV. Também são salvos exemplos visuais dos mapas, permitindo a inspeção qualitativa de acertos e erros. Uma rotina de agregação lê os arquivos de resumo das três sementes e calcula as médias e os desvios-padrão usados nas tabelas deste artigo.

Para reduzir o risco de confundir uma falha de instalação com um resultado científico, o notebook possui dois modos. O modo rápido usa uma época e uma fração pequena dos dados apenas para verificar o fluxo; o modo completo usa 15 épocas, o conjunto integral e 200 exemplos para Grad-CAM. Somente o segundo modo é utilizado nas tabelas e na discussão. A Tabela~\ref{tab:protocol} resume os principais parâmetros do protocolo.

\begin{table}[t]
\caption{Parâmetros principais da implementação.}
\label{tab:protocol}
\centering
\small
\begin{tabular}{ll}
\toprule
Parâmetro & Configuração \\
\midrule
Arquitetura & ResNet-18, pesos ImageNet \\
Classes & 37 raças de cães e gatos \\
Épocas no modo completo & 15 \\
Sementes & 42, 123 e 2026 \\
Exemplos para Grad-CAM & 200 por semente \\
\bottomrule
\end{tabular}
\end{table}


\section{Resultados}

\textcolor{red}{Seção a ser preenchida após os experimentos.}

\section{Discussão}

\textcolor{red}{Seção a ser preenchida após os resultados.}

\section{Conclusão}

\textcolor{red}{Seção a ser preenchida ao final.}

\bibliographystyle{IEEEtran}
\bibliography{references}

\end{document}
